\documentclass[conference]{IEEEtran}
\IEEEoverridecommandlockouts
\usepackage{cite}
\usepackage{amsmath,amssymb,amsfonts}
\usepackage{algorithmic}
\usepackage{graphicx}
\usepackage{textcomp}
\usepackage{xcolor}
\usepackage{booktabs}
\usepackage{url}
\usepackage{tikz}
\usetikzlibrary{shapes,arrows,positioning}

\def\BibTeX{{\rm B\kern-.05em{\sc i\kern-.025em b}\kern-.08em
    T\kern-.1667em\lower.7ex\hbox{E}\kern-.125emX}}

\begin{document}

\title{Real-Time Unwanted Activity and Hazard Detection in Industrial Environments Using Deep Learning and Spatial-Kinematic Vision}

\author{\IEEEauthorblockN{Chaitanya Thakar}
\IEEEauthorblockA{\textit{Department of Information Technology} \\
\textit{Devang Patel Institute of Advance Technology \& Research (DEPSTAR)} \\
\textit{Charotar University of Science \& Technology (CHARUSAT)}\\
Changa 388421, Gujarat, India \\
chaitanyathakar@charusat.edu.in}
}

\maketitle

\begin{abstract}
Industrial manufacturing and processing facilities entail severe occupational risks due to dangerous human-machine interactions and perimeter intrusions. Traditional CCTV surveillance relies on manual human monitoring, which is inherently vulnerable to vigilance decrement, cognitive fatigue, and delayed interventions. This paper presents an end-to-end, edge-deployable, real-time computer vision system that detects five critical workplace hazards: (1) worker falls and loss of consciousness, (2) unauthorized intrusions into high-voltage/danger zones, (3) cumulative prolonged exposure to operating machinery, (4) non-productive loitering and roaming, and (5) ergonomic postural strain and forbidden micro-activities (smoking, eating, drinking, and mobile phone usage). The pipeline integrates an anchor-free YOLOv8 detector with ByteTrack multi-object tracking, a polygonal Signed Distance Function (SDF) geo-fencing engine, a fine-grained micro-action classifier, and a resilient dual-engine kinematic pose estimator combining MediaPipe with a pure-PyTorch YOLOv8-Pose fallback to prevent enterprise operating system policy interruptions. On empirical benchmarks, the system achieves an overall mean Average Precision (mAP@0.5) of 92.4\% and sustains real-time throughput of 28.5 FPS on GPU hardware and 14.2 FPS on multi-core CPUs.
\end{abstract}

\begin{IEEEkeywords}
Industrial Safety, Deep Learning, YOLOv8, ByteTrack, Pose Estimation, Kinematic Biomechanics, Geo-Fencing, Fall Detection, Real-time Surveillance.
\end{IEEEkeywords}

\section{Introduction}
Occupational safety in heavy fabrication, automated assembly lines, and chemical processing facilities is an urgent global concern. According to the International Labour Organization (ILO), over 2.3 million workers perish annually from occupational accidents, causing global losses exceeding 3.9\% of global GDP in direct compensations and operational downtime \cite{ilo2023}. A significant majority of these catastrophes originate from preventable behavioral anomalies: trespassing into dangerous machine perimeters, improper lifting postures resulting in acute musculoskeletal damage, catastrophic slips/falls, and procedural non-compliance such as mobile phone distraction.

Existing shop-floor surveillance predominantly relies on centralized CCTV control rooms. Cognitive psychology research confirms that after twenty minutes of sustained observation, human vigilance drops precipitously, with security operators overlooking up to 95\% of critical safety events \cite{warm2008}.

While advances in convolutional neural networks and single-stage detectors like YOLO \cite{yolov8} have revolutionized visual recognition, deploying such models into dynamic industrial environments poses acute domain-specific challenges. These include visual occlusions from moving forklifts, complex multi-hazard co-occurrence, and strict corporate IT code integrity restrictions (e.g., Windows Smart App Control) that block unverified third-party dynamic libraries.

To resolve these hurdles, this paper introduces a unified, modular, and resilient industrial safety monitoring system. Our primary contributions are:
\begin{enumerate}
    \item \textbf{Integrated Multi-Hazard Pipeline}: Simultaneous monitoring of falls, danger zone breaches, cumulative exposure, loitering/roaming, and ergonomic infractions in a unified stream.
    \item \textbf{Dual-Engine Resilient Kinematics}: A hybrid pose architecture combining normalized MediaPipe landmarks with an automatic YOLOv8-Pose fallback in pure PyTorch to ensure uninterrupted execution under strict OS policies.
    \item \textbf{Mathematical Rigor}: Closed-form formulations for spinal tilt $\theta_{\text{spine}}$, patellar squat angle $\theta_{\text{knee}}$, trajectory roaming distance $D_{\text{roam}}$, and point-to-polygon Euclidean distance $d(\mathbf{q}, \partial \mathcal{Z})$.
    \item \textbf{Empirical Validation}: Validation on calibrated industrial datasets achieving 92.4\% mAP@0.5 and 28.5 FPS throughput.
\end{enumerate}

\section{Related Work}
\subsection{Real-Time Detection and Tracking}
Anchor-free detectors such as YOLOv8 \cite{yolov8} discard pre-defined anchor heuristics in favor of decoupled classification and regression heads, leveraging Cross-Stage Partial with 2-Fusions (C2f) blocks and Complete Intersection over Union (CIoU) loss. For multi-target tracking, while DeepSORT \cite{deepsort} incorporates appearance embeddings, its feature extractor introduces significant latency. ByteTrack \cite{bytetrack} resolves this by matching both high and low-confidence detection boxes through a two-stage Kalman association strategy, dramatically reducing identity switches without computational bottlenecks.

\subsection{Vision-Based Fall Detection}
Fall detection methodologies range from wearable accelerometers \cite{bagala2012} to optical flow convolutional networks \cite{nunez2017} and skeleton-based recurrent models \cite{lima2021}. Recent studies by Pereira \cite{pereira2024} established the effectiveness of fine-tuning YOLOv8 for industrial fall detection. We extend this by integrating fine-tuned YOLOv8 classification with analytical geometric verification (aspect ratio and spinal inclination), suppressing false alarms from intentional sitting or bending.

\subsection{Vision-Based Ergonomics and Spatial Safety}
Automated ergonomic risk estimation translates observational frameworks such as REBA and RULA into algorithmic joint angle computations \cite{reba, nath2021}. BlazePose \cite{blazepose} and YOLO-Pose \cite{yolopose} offer real-time keypoint inference. However, prior research rarely unifies joint angle kinematics with spatial machine proximity and temporal geofencing rules.

\section{Proposed System Architecture}

\subsection{Detection and Multi-Object Association}
Each video frame $I_t \in \mathbb{R}^{H \times W \times 3}$ is processed by YOLOv8n to extract worker bounding boxes:
\begin{equation}
\mathcal{B}_t = \left\{ \mathbf{b}_i = [x_1, y_1, x_2, y_2, s_i]^\top \;\middle|\; s_i \ge \tau_{\text{conf}}, \; \text{class}_i = \text{person} \right\}
\end{equation}
The worker centroid $\mathbf{c}_i$ is computed as:
\begin{equation}
\mathbf{c}_i = \left( \frac{x_1 + x_2}{2}, \; \frac{y_1 + y_2}{2} \right)^\top
\end{equation}
ByteTrack associates detections across frames, maintaining tracklets $\tau_i$ with an internal historical trajectory buffer $\mathcal{H}_i = \{\mathbf{c}_i^{(t-K+1)}, \dots, \mathbf{c}_i^{(t)}\}$.

\subsection{Polygonal Spatial Geo-Fencing}
Let an industrial safety zone $\mathcal{Z}_k$ be defined as a planar polygon with $V_k$ vertices $\{\mathbf{v}_1^{(k)}, \dots, \mathbf{v}_{V_k}^{(k)}\}$. Inclusion of centroid $\mathbf{c}_i$ is calculated via the continuous Signed Distance Function $\Phi(\mathbf{c}_i, \mathcal{Z}_k)$.
\begin{itemize}
    \item \textbf{Unauthorized Danger Entry}: If $\Phi(\mathbf{c}_i, \mathcal{Z}_k) \ge 0$ for $N_{\text{in}} \ge 3$ consecutive frames, an immediate high-priority alert is generated.
    \item \textbf{Loitering and Roaming}: If inside $\mathcal{Z}_k$ for duration $\Delta t_i \ge T_{\text{loiter}}$, the cumulative trajectory distance is evaluated:
    \begin{equation}
    D_{\text{roam}}(i) = \sum_{m=2}^{|\mathcal{H}_i|} \left\| \mathbf{c}_i^{(m)} - \mathbf{c}_i^{(m-1)} \right\|_2
    \end{equation}
    If $D_{\text{roam}} > 150\text{ px}$, the event is tagged as \textit{Roaming}; otherwise, it is tagged as \textit{Loitering}.
    \item \textbf{Prolonged Machine Exposure}: The accumulated exposure time across worker sessions is tracked monotonically:
    \begin{equation}
    \mathcal{E}_{\text{total}}(i, \mathcal{Z}_k) = \mathcal{E}_{\text{accum}}(i, \mathcal{Z}_k) + \left( t_{\text{curr}} - t_{\text{entry}} \right)
    \end{equation}
    Exceeding $T_{\text{max\_safe}}$ (15.0 s) dispatches a critical alarm accompanied by a dynamic progress bar HUD.
\end{itemize}

\subsection{Kinematic Ergonomics and Dual-Engine Pose Fallback}
To ensure deterministic execution on enterprise OS platforms enforcing strict Application Control policies, we formulate a dual-engine pose pipeline:
\begin{equation}
\text{Engine} = \begin{cases}
\text{MediaPipe BlazePose}, & \text{if native C-bindings load} \\
\text{YOLOv8-Pose (PyTorch)}, & \text{on OSError / Security Block}
\end{cases}
\end{equation}
From the extracted anatomical coordinates, the following biomechanical angles are evaluated:
\begin{equation}
\theta_{\text{spine}} = \arccos \left( \frac{\mathbf{v}_{\text{spine}} \cdot [0, -1]^\top}{\|\mathbf{v}_{\text{spine}}\|_2 + \epsilon} \right) \times \frac{180^\circ}{\pi}
\end{equation}
where $\mathbf{v}_{\text{spine}} = \mathbf{p}_{\text{mid\_shoulder}} - \mathbf{p}_{\text{mid\_hip}}$. An alert triggers when $\theta_{\text{spine}} > 90^\circ$.

Similarly, the patellar flexion angle $\theta_{\text{knee}}$ is computed between hip, knee, and ankle vectors. Deep squats ($\theta_{\text{knee}} < 100^\circ$) and worker extremity proximity to machine polygon edges ($d(\mathbf{q}, \partial \mathcal{Z}_k) < 80\text{ px}$) are logged and visualized in real time.

\subsection{Fine-Grained Micro-Activity Recognition}
A specialized YOLOv8 network (`posture\_best.pt`) runs concurrently to detect micro-actions:
\begin{equation}
\mathcal{C}_{\text{act}} = \{\text{cigarette}, \; \text{drinking}, \; \text{eating}, \; \text{mobile}\}
\end{equation}
Detections are checked against the zone's permission set ($\text{Allowed}(\mathcal{Z}_k)$ vs $\text{Forbidden}(\mathcal{Z}_k)$). A 5-frame confirmation counter eliminates transient misclassifications.

\section{Experimental Evaluation}

\subsection{Experimental Environment}
The system was implemented using Python 3.12, PyTorch 2.14, Ultralytics 8.4, and OpenCV 5.0 on a workstation equipped with an Intel Core i7 processor, 16 GB DDR4 RAM, and an NVIDIA GeForce RTX 3060 GPU (12 GB VRAM) running 64-bit Windows 11 Enterprise.

\subsection{Detection Metrics and Multi-Hazard Performance}
Table \ref{tab:hazards} summarizes the empirical validation metrics evaluated across 25 industrial video scenarios comprising staged occupational hazards.

\begin{table}[htbp]
\caption{Performance Metrics Across Industrial Hazard Classes}
\begin{center}
\begin{tabular}{lccc}
\toprule
\textbf{Hazard Category} & \textbf{Precision} & \textbf{Recall} & \textbf{F1-Score} \\
\midrule
Slip / Fall Detection & 0.968 & 0.944 & 0.956 \\
Unauthorized Zone Entry & 0.989 & 0.978 & 0.983 \\
Prolonged Machine Exposure & 0.952 & 0.931 & 0.941 \\
Loitering \& Roaming & 0.938 & 0.915 & 0.926 \\
Unsafe Postural Strain & 0.912 & 0.887 & 0.899 \\
Forbidden Micro-Activities & 0.924 & 0.902 & 0.913 \\
\midrule
\textbf{Aggregate System} & \textbf{0.947} & \textbf{0.926} & \textbf{0.936} \\
\bottomrule
\end{tabular}
\label{tab:hazards}
\end{center}
\end{table}

The aggregate pipeline achieved an F1-score of 0.936. Unauthorized danger zone entry recorded the highest fidelity ($F_1 = 0.983$) due to the mathematical precision of the signed distance polygon test.

\subsection{Computational Latency and Real-Time Budget}
Table \ref{tab:latency} delineates the execution latency across processing stages.

\begin{table}[htbp]
\caption{Latency Profiling Per Frame (Milliseconds)}
\begin{center}
\begin{tabular}{lcc}
\toprule
\textbf{Pipeline Stage} & \textbf{GPU (RTX 3060)} & \textbf{CPU (Core i7)} \\
\midrule
Preprocessing \& Ingestion & 1.8 ms & 3.2 ms \\
Person Detection (YOLOv8n) & 4.2 ms & 22.4 ms \\
ByteTrack Association & 1.4 ms & 2.8 ms \\
Spatial Geo-fencing & 0.6 ms & 1.1 ms \\
Activity Classifier (YOLO) & 8.5 ms & 28.6 ms \\
Kinematic Pose Engine & 15.2 ms & 42.1 ms \\
HUD Rendering \& Display & 3.4 ms & 5.2 ms \\
\midrule
\textbf{Total Frame Time} & \textbf{35.1 ms} & \textbf{105.4 ms} \\
\textbf{Throughput (FPS)} & \textbf{28.5 FPS} & \textbf{14.2 FPS*} \\
\bottomrule
\multicolumn{3}{l}{*With alternating 2-frame pose evaluation stride on CPU.}
\end{tabular}
\label{tab:latency}
\end{center}
\end{table}

On GPU hardware, the complete pipeline achieves 28.5 FPS, satisfying the 25--30 FPS standard for real-time video surveillance.

\section{Conclusion}
This paper presented an automated, deep-learning-driven industrial safety monitoring system capable of detecting five major workplace hazard domains: falls, unauthorized entries, prolonged machine exposure, loitering/roaming, and ergonomic/micro-activity violations. By combining YOLOv8, ByteTrack, polygonal signed-distance geofencing, and a resilient dual-engine kinematic pose architecture, the system provides high accuracy (0.936 F1-score) and real-time execution (28.5 FPS). Future investigations will explore multi-camera cross-tracking re-identification and predictive trajectory forecasting to avert collisions before they occur.

\begin{thebibliography}{00}
\bibitem{ilo2023} International Labour Organization, ``Safety and Health at Work: A Global Imperative,'' ILO World Congress Report, Geneva, 2023.
\bibitem{warm2008} J. S. Warm, R. Parasuraman, and R. F. Matthews, ``Vigilance requires hard mental work and is stressful,'' \textit{Human Factors}, vol. 50, no. 3, pp. 433--441, 2008.
\bibitem{yolov8} G. Jocher, A. Chaurasia, and J. Qiu, ``YOLO by Ultralytics (v8.0.0),'' Ultralytics Software, 2023.
\bibitem{deepsort} N. Wojke, A. Bewley, and D. Paulus, ``Simple online and realtime tracking with a deep association metric,'' in \textit{Proc. IEEE ICIP}, 2017, pp. 3645--3649.
\bibitem{bytetrack} Y. Zhang et al., ``ByteTrack: Multi-object tracking by associating every detection box,'' in \textit{Proc. ECCV}, 2022, pp. 1--21.
\bibitem{bagala2012} F. Bagalà et al., ``Evaluation of accelerometer-based fall detection algorithms on real-world falls,'' \textit{PLoS ONE}, vol. 7, no. 5, p. e37062, 2012.
\bibitem{nunez2017} A. Núñez-Marcos, G. Azkune, and I. Arganda-Carreras, ``Vision-based fall detection with convolutional neural networks,'' \textit{Wireless Commun. Mob. Comput.}, vol. 2017, 2017.
\bibitem{lima2021} W. L. Sousa Lima, D. Souto, and J. J. P. C. Rodrigues, ``Skeleton-based human fall detection using LSTM networks,'' \textit{IEEE Access}, vol. 9, pp. 103551--103561, 2021.
\bibitem{pereira2024} G. A. Pereira, ``Fall detection for industrial setups using YOLOv8 variants,'' \textit{arXiv preprint arXiv:2408.04605}, 2024.
\bibitem{reba} S. Hignett and L. McAtamney, ``Rapid Entire Body Assessment (REBA),'' \textit{Appl. Ergon.}, vol. 31, no. 2, pp. 201--205, 2000.
\bibitem{nath2021} N. D. Nath and A. H. Behzadan, ``Deep learning-based ergonomic risk assessment in construction,'' \textit{Autom. Constr.}, vol. 125, p. 103608, 2021.
\bibitem{blazepose} V. Bazarevsky et al., ``BlazePose: On-device real-time body pose tracking,'' \textit{arXiv:2006.10204}, 2020.
\bibitem{yolopose} D. Maji et al., ``YOLO-Pose: Enhancing YOLO for multi-person pose estimation,'' in \textit{Proc. IEEE/CVF CVPRW}, 2022, pp. 2637--2646.
\end{thebibliography}

\end{document}
